\documentclass[roman, 12pt]{beamer}

%%% Работа с русским языком
\usepackage{mathtext} 				% русские буквы в формулах
\usepackage[T2A]{fontenc}			% кодировка
\usepackage[utf8]{inputenc}			% кодировка исходного текста
\usepackage[english,russian]{babel}	% локализация и переносы

%%% Дополнительная работа с математикой
\usepackage{amsmath, amsthm, amssymb, eucal}

\usepackage{geometry} % Простой способ задавать поля
\geometry{top=20mm}
\geometry{bottom=20mm}
\geometry{left=25mm}
\geometry{right=25mm}

\usepackage{graphicx}
\graphicspath{{D:/latex-tasks/lab-6/images}{images/}} % Папка с картинками
\setlength\fboxsep{3pt} % Отступ рамки \fbox{} от рисунка
\setlength\fboxrule{1pt} % Толщина линий рамки fbox{}
%%% Выбор стиля страниц
\pagestyle{plain}

\usetheme{Antibes}
%\usecolortheme{}%

%%% Заголовок
\title{Топ 5 смешариков}
\date{17 мая 2022 г.}
\author{Шпак А.В.}

%Привет, друг :)%

\begin{document} % Конец преамбулы, начало текста.

	\begin{frame}
		\titlepage
	\end{frame}
	
	\begin{frame}
		\frametitle{Пятое место.}
		
		\begin {columns}
		\column{0.7\textwidth}
		\textcolor{brown}{Копатыч - садовод и огородник. Простой и приятный смешарик. Увлекался диско. Находиться в его компании - одно удовольствие!}
		\column{0.7\textwidth}
		\begin{figure}
			\includegraphics[width=0.7\linewidth]{images/kopatich}
		\end{figure}
		\end {columns}
		
	\end {frame}
	
	\begin{frame}
		\frametitle{Четвертое место.}
		
		\begin {columns}
		\column{0.7\textwidth}
		\begin{figure}
			\includegraphics[width=0.7\linewidth]{images/losyash}
		\end{figure}
		\column{0.7\textwidth}
		\textcolor{yellow}{Лосяш - умнейший из смешариков. Проводит опыты и любит бабочек. Дружит с Копатычем}
		\end {columns}		
	\end {frame}
		
	\begin{frame}
		\frametitle{Третье место.}
		
		\begin {columns}
		\column{0.7\textwidth}
		\textcolor{blue}{Каркарыч, кажется, самый адекватный из всех смешариков. Реалист, мечтатель, он интеллегентен и вежлив. Обладает богатым жизненным опытом и обожает игру бильярд.}
		\column{0.7\textwidth}
		\begin{figure}
			\includegraphics[width=0.7\linewidth]{images/karkarich}
		\end{figure}
		\end {columns}		
		
	\end {frame}
			
	\begin{frame}
		\frametitle{Второе место.}
		
		\begin {columns}
		\column{0.7\textwidth}
		\begin{figure}
			\includegraphics[width=0.7\linewidth]{images/ezhik}
		\end{figure}
		\column{0.7\textwidth}
		\textcolor{green}{Ёжик - самый лучший друг Кроша и самый хороший смешарик. Он умный и понимающий, чуткий. Любит кактусы.}
		\end {columns}
		
		\end {frame}
	
	\begin{frame}
		\frametitle{Первое место.}
		
		\begin{block}{Победитель топа смешариков - Совунья!!}
		\end{block}
		
		\begin {columns}
		\column{0.7\textwidth}
		\textcolor{violet}{Совунья - самый душевный, родной и добрый смешарик. Она всегда помогала другим, лечила других и не бросала в беде!}
		\column{0.7\textwidth}
		\begin{figure}
			\includegraphics[width=0.7\linewidth]{images/sovunia}
		\end{figure}
		\end {columns}
	
		\end {frame}
	
	\begin{frame}
		\frametitle{Конец!}
		\end {frame}
	
\end{document} % Конец текста.